% ============ Appendix to §6. Secondary representations of a module (原书 pp.42-44) ============
\section*{Appendix to \S\ref{sec:6}. Secondary representations of a module}
\phantomsection
\addcontentsline{toc}{section}{Appendix to \S\ref{sec:6}. Secondary representations of a module}

I.G. Macdonald~\cite{MacD1} has developed the theory of \emph{attached prime ideals}\index{attached prime ideal}
and \emph{secondary representations}\index{secondary representation} of a module, which is in a certain sense
dual to the theory of associated prime ideals and primary decompositions.
This theory was successfully applied to the theory of local cohomology
by him and R.Y. Sharp (Macdonald \& Sharp~\cite{MS1}, Sharp~\cite{Shar7}).

Let $A$ be a commutative ring. An $A$-module $M$ is said to be \emph{secondary}\index{secondary module}
if $M \neq 0$ and, for each $a \in A$, the endomorphism $\varphi_a\colon M \to M$ defined
by $\varphi_a(m) = am$ (for $m \in M$) is either surjective or nilpotent. If $M$ is secondary,
then $P = \sqrt{\ann M}$ is a prime ideal, and $M$ is said to be \emph{$P$-secondary}. Any
non-zero quotient of a $P$-secondary module is $P$-secondary.

\begin{example}\label{ex:6.5}
If $A$ is an integral domain, its quotient field $K$ is a $(0)$-secondary
$A$-module.
\end{example}

\begin{example}\label{ex:6.6}
Let $W = \Z[p^{-1}]$, where $p$ is a prime number, and consider the
Artinian $\Z$-module $W/\Z$ (see \S\ref{sec:3}). This is also a $(0)$-secondary $\Z$-module.
\end{example}

\begin{example}\label{ex:6.7}
If $A$ is a local ring with maximal ideal $P$ and if every element of
$p$ is nilpotent, then $A$ itself is a $P$-secondary $A$-module.
\end{example}

\begin{example}\label{ex:6.8}
If $P$ is a maximal ideal of $A$, then $A/P^n$ is a $P$-secondary $A$-
module for every $n > 0$.
\end{example}

A \emph{secondary representation} of an $A$-module $M$ is an expression of $M$ as a
finite sum of secondary submodules:
\[
  (*) \quad M = N_1 + \cdots + N_r.
\]
The representation is \emph{minimal}\index{minimal secondary representation} if (1) the prime ideals $P_i := \sqrt{\ann N_i}$ are all
distinct, and (2) none of the $N_i$ is redundant. It is easy to see that the sum of
two $P$-secondary submodules is again $P$-secondary, hence if $M$ has a
secondary representation then it has a minimal one.

A prime ideal $P$ is called an \emph{attached prime ideal} of $M$ if $M$ has a $P$-
secondary quotient. The set of the attached prime ideals of $M$ is denoted by
$\Att(M)$.

\begin{theorem}\label{thm:6.9}
If $(*)$ is a minimal secondary representation of $M$ and $P_i =
\sqrt{\ann N_i}$, then $\Att(M) = \{P_1, \ldots, P_n\}$.
\end{theorem}

\begin{proof}
Since $M/(N_1 + \cdots + N_{i-1} + N_{i+1} + \cdots + N_n)$ is a non-zero quotient
of $N_i$, it is a $P_i$-secondary module. Thus $\{P_1, \ldots, P_n\} \subset \Att(M)$. Conver-
sely, let $P \in \Att(M)$ and let $W$ be a $P$-secondary quotient of $M$. Then $W =
W_1 + \cdots + W_n$, where $W_i$ is the image of $N_i$ in $W$. From this we obtain a
minimal secondary representation $W = W'_1 + \cdots + W'_m$, and then
$\Att(W) \supset \{P'_1, \ldots, P'_m\}$. On the other hand $\Att(W) = \{P\}$ since $W$ is
$P$-secondary. Therefore $P = P_i$ for some $i$.
\end{proof}

\begin{theorem}\label{thm:6.10}
If $0 \to M' \to M \to M'' \to 0$ is an exact sequence of
$A$-modules, then $\Att(M'') \subset \Att(M) \subset \Att(M') \cup \Att(M'')$.
\end{theorem}

\begin{proof}
The first inclusion is trivial from the definition. For the second, let
$P \in \Att(M)$ and let $N$ be a submodule such that $M/N$ is $P$-secondary. If
$M' + N = M$ then $M/N$ is a non-trivial quotient of $M'$, hence $P \in \Att(M')$.
If $M' + N \neq M$ then $M/(M' + N)$ is a non-trivial quotient of $M''$ as well
as of $M/N$, hence $M''$ has a $P$-secondary quotient and $P \in \Att(M'')$.
\end{proof}

An $A$-module $M$ is said to be \emph{sum-irreducible}\index{sum-irreducible module} if it is neither zero nor the
sum of two proper submodules.

\begin{lemma*}\label{lem:6.1}
If $M$ is Artinian and sum-irreducible, then it is secondary.
\end{lemma*}

\begin{proof}
Suppose $M$ is not secondary. Then there is $a \in A$ such that $M \neq aM$
and $a^nM \neq 0$ for all $n > 0$. Since $M$ is Artinian, we have $a^nM = a^{n+1}M$ for
some $n$. Set $K = \{x \in M \mid a^nx = 0\}$. Then it is immediate that $M = K + aM$,
and so $M$ is not sum-irreducible.
\end{proof}

\begin{theorem}\label{thm:6.11}
If $M$ is Artinian, then it has a secondary representation.
\end{theorem}

\begin{proof}
Similar to the proof of Theorem~\ref{thm:6.8}, (iii).
\end{proof}

The class of modules which have secondary representations is larger
than that of Artinian modules. Sharp [8] proved that an injective module
over a Noetherian ring has a secondary representation.

% ============ 习题 ============
\section*{Exercises to Appendix to \S\ref{sec:6}}

\begin{exercise}\label{ex:6.8-ex}
An $A$-module $M$ is \emph{coprimary}\index{coprimary module} if $\Ass(M)$ has just one element. Show that a
finite module $M \neq 0$ over a Noetherian ring $A$ is coprimary if and only if
the following condition is satisfied: for every $a \in A$, the endomorphism
$a\colon M \to M$ is either injective or nilpotent. In this case $\Ass M = \{P\}$,
where $P = \sqrt{\ann M}$.
\end{exercise}

\begin{exercise}\label{ex:6.9-ex}
Show that if $M$ is an $A$-module of finite length then $M$ is coprimary if and
only if it is secondary. Show also that such a module $M$ is a direct sum of
secondary modules belonging to maximal ideals, and $\Ass(M) = \Att(M)$.
\end{exercise}
